% !TEX program = lualatex
% ---------------------------------------------------------------------------------------------
%  Readback note 3: The five final theorems, and what one must trust
%
%  Build:  python3 tools/rb.py note <slug>
%
%  One unit of a Lean artifact, rewritten as ordinary mathematics. The prose is hand-written. Lean names,
%  line ranges and counts come from data.tex; quoted Lean statements are copied from the artifact by the build.
%
%  Rules of the text: standard terminology; every symbol is explained where it first appears and listed again
%  in the notation at the end; no letter has two meanings; a proof is complete or says what it takes for granted;
%  Lean names appear only in the LEAN lines and in the concordance.
%
%  The signature rule (the build enforces it): a function, relation or predicate that the artifact defines is
%  written only through its notation macro, which the build makes from the spec (`notation`) with one slot for
%  each explicit argument of the Lean definition. No argument is fixed and dropped, also not "for brevity" and
%  not for one section. Every argument and every hypothesis of a stated declaration stands inside its
%  statement, or in a Setting that the statement cites by number; only implicit and instance arguments (the
%  ambient objects) may stand in the conventions.
% ---------------------------------------------------------------------------------------------
\documentclass[11pt,a4paper]{article}
\ifdefined\pdfvariable\pdfvariable suppressoptionalinfo \numexpr 1+2+4+8\relax\fi
\usepackage[a4paper,hmargin=35mm,top=22mm,bottom=27mm,footskip=12mm,nohead]{geometry}

\usepackage{fontspec}
\directlua{
  luaotfload.add_fallback("gfallback", {"STIXTwoMath-Regular.otf:mode=node;", "DejaVuSans.ttf:mode=node;"})
  luaotfload.add_fallback("gmonofallback", {"DejaVuSansMono.ttf:mode=node;", "STIXTwoMath-Regular.otf:mode=node;"})
}
\setmainfont{NewCM10}[Extension=.otf,UprightFont=*-Book,ItalicFont=*-BookItalic,
  BoldFont=*-Bold,BoldItalicFont=*-BoldItalic,SmallCapsFont=*-Book,SmallCapsFeatures={Letters=SmallCaps},
  RawFeature={fallback=gfallback}]
\setsansfont{NewCMSans10}[Extension=.otf,UprightFont=*-Book,ItalicFont=*-BookOblique,
  BoldFont=*-Bold,BoldItalicFont=*-BoldOblique,RawFeature={fallback=gfallback}]
\setmonofont{JetBrainsMono}[Extension=.otf,UprightFont=*-Regular,ItalicFont=*-Italic,
  BoldFont=*-SemiBold,BoldItalicFont=*-SemiBoldItalic,Scale=MatchLowercase,
  RawFeature={-calt;-liga;fallback=gmonofallback}]
\usepackage{amsmath}
\usepackage{amsthm}
\usepackage[math-style=TeX]{unicode-math}
\setmathfont{NewCMMath-Book.otf}
\emergencystretch=2.5em\relax

\usepackage{xcolor}
\usepackage{tikz}
\usetikzlibrary{calc,arrows.meta,positioning,backgrounds}
\usepackage[protrusion=true,expansion=true]{microtype}
\usepackage{enumitem}
\usepackage{needspace}
\usepackage{booktabs}
\usepackage{array}
\usepackage{longtable}
\usepackage{fancyhdr}
\usepackage{lastpage}
\usepackage{multicol}
\usepackage[hidelinks,pdfusetitle]{hyperref}

\definecolor{ink}{HTML}{0F172A}
\definecolor{inksoft}{HTML}{475569}
\definecolor{inkfaint}{HTML}{94A3B8}
\definecolor{rulegrey}{HTML}{CBD5E1}
\definecolor{arith}{HTML}{C2410C}
\definecolor{est}{HTML}{7C3AED}
\definecolor{ana}{HTML}{0891B2}
\definecolor{ok}{HTML}{047857}
\definecolor{leancol}{HTML}{3730A3}

\title{The five final theorems, and what one must trust}
\author{grh-compact}
\date{}

% ------------------------------------------------------------------------------------- macros
\newcommand{\leanfont}{\ttfamily\fontsize{8.6}{10.5}\selectfont\color{leancol}}
\newcommand{\lean}[1]{\csname LN@#1\endcsname}
\newcommand{\fillin}[1]{{\color{arith}\itshape[#1]}}  % a place of the template that is not written yet; the build reports every one that is left
\newcommand{\leanquote}[1]{\csname LQ@#1\endcsname}
% generated by rb.py note; do not edit
\newcommand{\ConcDecls}{12}
\newcommand{\Pin}{72fead10}
\newcommand{\PinText}{that of commit {\leanfont 72fead10}}
\newcommand{\ArtifactName}{grh-compact}
\expandafter\newcommand\csname LN@thmRH\endcsname{\href{https://github.com/lsernau/grh-lean/blob/72fead102f663d5c7a407964a13dc0cf9d977207/grh-compact/src/Main.lean\#L4-L4}{{\leanfont riemannHypothesis}}}
\expandafter\newcommand\csname LQ@thmRH\endcsname{\leanquoteline{0}{theorem riemannHypothesis : RiemannHypothesis}}
\expandafter\newcommand\csname LN@thmZeta\endcsname{\href{https://github.com/lsernau/grh-lean/blob/72fead102f663d5c7a407964a13dc0cf9d977207/grh-compact/src/Main.lean\#L6-L6}{{\leanfont riemannZeta\_eq\_zero\_re}}}
\expandafter\newcommand\csname LQ@thmZeta\endcsname{\leanquoteline{0}{theorem riemannZeta\_eq\_zero\_re : RiemannZetaZerosOnLine}}
\expandafter\newcommand\csname LN@thmHalf\endcsname{\href{https://github.com/lsernau/grh-lean/blob/72fead102f663d5c7a407964a13dc0cf9d977207/grh-compact/src/Main.lean\#L8-L8}{{\leanfont LFunction\_ne\_zero\_of\_half\_lt\_re}}}
\expandafter\newcommand\csname LQ@thmHalf\endcsname{\leanquoteline{0}{theorem LFunction\_ne\_zero\_of\_half\_lt\_re : DirichletLNonvanishingRightOfHalf}}
\expandafter\newcommand\csname LN@thmStrip\endcsname{\href{https://github.com/lsernau/grh-lean/blob/72fead102f663d5c7a407964a13dc0cf9d977207/grh-compact/src/Main.lean\#L10-L10}{{\leanfont LFunction\_eq\_zero\_re}}}
\expandafter\newcommand\csname LQ@thmStrip\endcsname{\leanquoteline{0}{theorem LFunction\_eq\_zero\_re : DirichletLZerosOnLine}}
\expandafter\newcommand\csname LN@thmPrim\endcsname{\href{https://github.com/lsernau/grh-lean/blob/72fead102f663d5c7a407964a13dc0cf9d977207/grh-compact/src/Main.lean\#L12-L13}{{\leanfont LFunction\_eq\_zero\_re\_of\_isPrimitive}}}
\expandafter\newcommand\csname LQ@thmPrim\endcsname{\leanquoteline{0}{theorem LFunction\_eq\_zero\_re\_of\_isPrimitive : PrimitiveDirichletLZerosOnLine}}
\expandafter\newcommand\csname LN@grh\endcsname{\href{https://github.com/lsernau/grh-lean/blob/72fead102f663d5c7a407964a13dc0cf9d977207/grh-compact/src/New/Main.lean\#L174-L190}{{\leanfont New.\allowbreak{}grh}}}
\expandafter\newcommand\csname LQ@grh\endcsname{\leanquoteline{0}{lemma New.grh : DirichletLNonvanishingRightOfHalf}}
\expandafter\newcommand\csname LN@zerosOnLine\endcsname{\href{https://github.com/lsernau/grh-lean/blob/72fead102f663d5c7a407964a13dc0cf9d977207/grh-compact/src/New/Main.lean\#L194-L196}{{\leanfont New.\allowbreak{}zeros\_on\_line}}}
\expandafter\newcommand\csname LQ@zerosOnLine\endcsname{\leanquoteline{0}{lemma New.zeros\_on\_line (h : DirichletLNonvanishingRightOfHalf) : DirichletLZerosOnLine}}
\expandafter\newcommand\csname LN@zetaZerosOnLine\endcsname{\href{https://github.com/lsernau/grh-lean/blob/72fead102f663d5c7a407964a13dc0cf9d977207/grh-compact/src/New/Main.lean\#L200-L204}{{\leanfont New.\allowbreak{}zeta\_zeros\_on\_line}}}
\expandafter\newcommand\csname LQ@zetaZerosOnLine\endcsname{\leanquoteline{0}{lemma New.zeta\_zeros\_on\_line (h : DirichletLNonvanishingRightOfHalf) :}\leanquoteline{4}{RiemannZetaZerosOnLine}}
\expandafter\newcommand\csname LN@rh\endcsname{\href{https://github.com/lsernau/grh-lean/blob/72fead102f663d5c7a407964a13dc0cf9d977207/grh-compact/src/Statements.lean\#L3-L4}{{\leanfont RiemannZetaZerosOnLine}}}
\expandafter\newcommand\csname LN@halfPlane\endcsname{\href{https://github.com/lsernau/grh-lean/blob/72fead102f663d5c7a407964a13dc0cf9d977207/grh-compact/src/Statements.lean\#L6-L8}{{\leanfont DirichletLNonvanishingRightOfHalf}}}
\expandafter\newcommand\csname LN@strip\endcsname{\href{https://github.com/lsernau/grh-lean/blob/72fead102f663d5c7a407964a13dc0cf9d977207/grh-compact/src/Statements.lean\#L10-L12}{{\leanfont DirichletLZerosOnLine}}}
\expandafter\newcommand\csname LN@stripPrim\endcsname{\href{https://github.com/lsernau/grh-lean/blob/72fead102f663d5c7a407964a13dc0cf9d977207/grh-compact/src/Statements.lean\#L14-L16}{{\leanfont PrimitiveDirichletLZerosOnLine}}}
\newcommand{\ConcordanceRows}{\crow{thmRH}{\href{https://github.com/lsernau/grh-lean/blob/72fead102f663d5c7a407964a13dc0cf9d977207/grh-compact/src/Main.lean\#L4-L4}{{\leanfont riemannHypothesis}}}{Main.lean}{4--4}
\crow{thmZeta}{\href{https://github.com/lsernau/grh-lean/blob/72fead102f663d5c7a407964a13dc0cf9d977207/grh-compact/src/Main.lean\#L6-L6}{{\leanfont riemannZeta\_eq\_zero\_re}}}{Main.lean}{6--6}
\crow{thmHalf}{\href{https://github.com/lsernau/grh-lean/blob/72fead102f663d5c7a407964a13dc0cf9d977207/grh-compact/src/Main.lean\#L8-L8}{{\leanfont LFunction\_ne\_zero\_of\_half\_lt\_re}}}{Main.lean}{8--8}
\crow{thmStrip}{\href{https://github.com/lsernau/grh-lean/blob/72fead102f663d5c7a407964a13dc0cf9d977207/grh-compact/src/Main.lean\#L10-L10}{{\leanfont LFunction\_eq\_zero\_re}}}{Main.lean}{10--10}
\crow{thmPrim}{\href{https://github.com/lsernau/grh-lean/blob/72fead102f663d5c7a407964a13dc0cf9d977207/grh-compact/src/Main.lean\#L12-L13}{{\leanfont LFunction\_eq\_zero\_re\_of\_isPrimitive}}}{Main.lean}{12--13}
\crow{grh}{\href{https://github.com/lsernau/grh-lean/blob/72fead102f663d5c7a407964a13dc0cf9d977207/grh-compact/src/New/Main.lean\#L174-L190}{{\leanfont New.\allowbreak{}grh}}}{Main.lean}{174--190}
\crow{zerosOnLine}{\href{https://github.com/lsernau/grh-lean/blob/72fead102f663d5c7a407964a13dc0cf9d977207/grh-compact/src/New/Main.lean\#L194-L196}{{\leanfont New.\allowbreak{}zeros\_on\_line}}}{Main.lean}{194--196}
\crow{zetaZerosOnLine}{\href{https://github.com/lsernau/grh-lean/blob/72fead102f663d5c7a407964a13dc0cf9d977207/grh-compact/src/New/Main.lean\#L200-L204}{{\leanfont New.\allowbreak{}zeta\_zeros\_on\_line}}}{Main.lean}{200--204}
\crow{rh}{\href{https://github.com/lsernau/grh-lean/blob/72fead102f663d5c7a407964a13dc0cf9d977207/grh-compact/src/Statements.lean\#L3-L4}{{\leanfont RiemannZetaZerosOnLine}}}{Statements.lean}{3--4}
\crow{halfPlane}{\href{https://github.com/lsernau/grh-lean/blob/72fead102f663d5c7a407964a13dc0cf9d977207/grh-compact/src/Statements.lean\#L6-L8}{{\leanfont DirichletLNonvanishingRightOfHalf}}}{Statements.lean}{6--8}
\crow{strip}{\href{https://github.com/lsernau/grh-lean/blob/72fead102f663d5c7a407964a13dc0cf9d977207/grh-compact/src/Statements.lean\#L10-L12}{{\leanfont DirichletLZerosOnLine}}}{Statements.lean}{10--12}
\crow{stripPrim}{\href{https://github.com/lsernau/grh-lean/blob/72fead102f663d5c7a407964a13dc0cf9d977207/grh-compact/src/Statements.lean\#L14-L16}{{\leanfont PrimitiveDirichletLZerosOnLine}}}{Statements.lean}{14--16}}
% the notations of the defined functions: one slot for each explicit argument of the Lean declaration
\newcommand{\UnitMembers}{5}
\newcommand{\UnitInNote}{5}
\newcommand{\UnitProofs}{0}
\newcommand{\UnitSteps}{0}
\newcommand{\UnitAllSteps}{0}


\linespread{1.08}                       % 11 pt on 14.7 pt
\setlength{\parindent}{0pt}
\setlength{\parskip}{5.5pt plus 1.5pt minus 0.5pt}
\setlength{\abovedisplayskip}{7.5pt plus 2pt minus 3pt}
\setlength{\belowdisplayskip}{7.5pt plus 2pt minus 3pt}
\setlength{\abovedisplayshortskip}{3pt plus 2pt}
\setlength{\belowdisplayshortskip}{6pt plus 2pt minus 2pt}
\emergencystretch=1.5em
\clubpenalty=10000 \widowpenalty=10000 \displaywidowpenalty=10000
\relpenalty=7000 \binoppenalty=9000
\raggedbottom
\pagestyle{fancy}\fancyhf{}\renewcommand{\headrulewidth}{0pt}
\fancyfoot[C]{\sffamily\fontsize{7.6}{9}\selectfont\color{inkfaint}\thepage\,/\,\pageref*{LastPage}}
\setlist{topsep=2pt,partopsep=0pt,parsep=0pt,itemsep=4pt,leftmargin=2.1em,labelsep=0.55em}

\newcommand{\Q}{\mathbb{Q}}
\newcommand{\Z}{\mathbb{Z}}
\newcommand{\F}{\mathbb{F}}
\newcommand{\len}{\operatorname{len}}
\newcommand{\cok}{\operatorname{cok}}
\newcommand{\End}{\operatorname{End}}
\newcommand{\Hom}{\operatorname{Hom}}
\newcommand{\Gal}{\operatorname{Gal}}
\newcommand{\diag}{\operatorname{diag}}
\newcommand{\kw}[1]{\emph{#1}}
\newcommand{\m}[1]{\mbox{$#1$}}        % a short formula that must not be broken across lines
\newcommand{\head}[1]{\par\addvspace{17pt plus 4pt minus 2pt}\Needspace{6\baselineskip}%
  \noindent{\fontsize{13.2}{16}\selectfont\bfseries\color{ink}#1}\par\nopagebreak\vspace{-1pt}}
\newcommand{\sect}[2]{\refstepcounter{section}\label{#2}\par\addvspace{19pt plus 4pt minus 2pt}\Needspace{7\baselineskip}%
  \noindent\llap{{\fontsize{13.2}{16}\selectfont\bfseries\color{inkfaint}\thesection}\hspace{9pt}}%
  {\fontsize{13.2}{16}\selectfont\bfseries\color{ink}#1}\par\nopagebreak\vspace{-1pt}}
\newcommand{\subhead}[1]{\par\addvspace{9pt plus 2pt}\Needspace{4\baselineskip}\noindent\textbf{#1}\enspace\ignorespaces}
% where a statement lives in the Lean archive
\newcommand{\leantag}{{\sffamily\bfseries\fontsize{6.3}{8}\selectfont\color{leancol}\addfontfeatures{LetterSpace=7}LEAN}}
\newcommand{\sep}{\unskip\nobreak\hspace{5pt}{\color{inkfaint}\textperiodcentered}\hspace{5pt plus 2pt}\ignorespaces}
\newcommand{\leanline}[1]{\par\nopagebreak\vspace{-2.5pt}{\fontsize{8.9}{12.4}\selectfont\color{inksoft}\raggedright
  \noindent\llap{\leantag\hspace{9pt}}#1\par}\vspace{1pt}}
% a displayed statement in a tinted box: #1 colour, #2 text
\newcommand{\statement}[2]{\par\addvspace{7pt}\begingroup\setlength{\fboxsep}{8.5pt}\setlength{\fboxrule}{0.6pt}%
  \noindent\fcolorbox{#1!30!white}{#1!3!white}{\parbox{\dimexpr\linewidth-2\fboxsep-2\fboxrule\relax}{#2}}%
  \endgroup\par\addvspace{7pt}}
\newcommand{\nt}[2]{\par\noindent\hangindent=1.1em\hangafter=1 #1\enspace{\color{inksoft}#2}\par\vspace{1pt}}
\newcommand{\legdot}[1]{\tikz[baseline=-0.62ex]\fill[#1] (0,0) circle (2.5pt);\hspace{3pt}}
\newcommand{\crow}[4]{\csname CS@#1\endcsname & #2 & \textcolor{inksoft}{#3} & \textcolor{inksoft}{#4}\\}
\newcommand{\leanquoteline}[2]{\par\noindent\hangindent=\dimexpr#1\fontdimen2\font+2.2em\relax\hangafter=1
  \hspace*{#1\fontdimen2\font}#2}
\newcommand{\quoted}[1]{\par\vspace{3pt}{\ttfamily\fontsize{7.5}{10.4}\selectfont\color{leancol}\raggedright
  \setlength{\parskip}{0pt}\leanquote{#1}\par}\vspace{4pt}}
\newcommand{\bx}[3]{{\sffamily\bfseries\fontsize{7.3}{9.2}\selectfont\color{#1!82!black}#2}\\[1.5pt]#3}
\newcommand{\cs}[2]{\expandafter\newcommand\csname CS@#1\endcsname{#2}}

% theorems
\newtheoremstyle{plainx}{7pt plus 2pt}{7pt plus 2pt}{\itshape}{}{\bfseries}{.}{0.5em}{}
\newtheoremstyle{uprightx}{7pt plus 2pt}{7pt plus 2pt}{}{}{\bfseries}{.}{0.5em}{}
\theoremstyle{plainx}
\newtheorem{theorem}{Theorem}[section]
\newtheorem{lemma}[theorem]{Lemma}
\newtheorem{proposition}[theorem]{Proposition}
\newtheorem{corollary}[theorem]{Corollary}
\theoremstyle{uprightx}
\newtheorem{definition}[theorem]{Definition}
\newtheorem{setting}[theorem]{Setting}          % hypotheses shared by several statements; each of them cites it by number
\newtheorem{remark}[theorem]{Remark}
\newtheorem{example}[theorem]{Example}
\renewcommand{\qedsymbol}{{\setlength{\fboxsep}{0pt}\setlength{\fboxrule}{0.5pt}\fbox{\rule{0pt}{0.46em}\hspace{0.46em}}}}
% what holds throughout the note: the ambient objects only (implicit and instance arguments), never a hypothesis
\newenvironment{conventions}{\par\addvspace{5pt}\noindent\textbf{Conventions.}\enspace\ignorespaces}{\par\addvspace{3pt}}
\makeatletter
\renewenvironment{proof}[1][Proof]{\par\pushQED{\qed}\normalfont\topsep0pt\partopsep0pt
  \trivlist\item[\hskip\labelsep\itshape #1.]\ignorespaces}{\popQED\endtrivlist\@endpefalse}
\makeatother

% where each declaration of the concordance is treated: \cs{key}{place}, the place naming a \label or a lettered item
\cs{thmRH}{Theorem 1}
\cs{thmZeta}{Theorem 2}
\cs{thmHalf}{Theorem 3}
\cs{thmStrip}{Theorem 4}
\cs{thmPrim}{Theorem 5}
\cs{grh}{Input 1}
\cs{zerosOnLine}{Input 2}
\cs{zetaZerosOnLine}{Input 3}
\cs{rh}{Definition~\ref{d:statements}}
\cs{halfPlane}{Definition~\ref{d:statements}}
\cs{strip}{Definition~\ref{d:statements}}
\cs{stripPrim}{Definition~\ref{d:statements}}
\newcommand{\CC}{\mathbb{C}}
\newcommand{\ZZ}{\mathbb{Z}}
\newcommand{\re}{\operatorname{Re}}

\begin{document}
\color{ink}

{\raggedright
{\sffamily\bfseries\fontsize{7.2}{9}\selectfont\color{leancol}\addfontfeatures{LetterSpace=8}READBACK NOTE 3\par}
\vspace{5pt}
{\fontsize{21.5}{25.5}\selectfont\bfseries The five final theorems, and what one must trust\par}
\vspace{5pt}
{\fontsize{12.4}{16}\selectfont\itshape\color{inksoft} One unit of the Lean artifact \ArtifactName, read back into
mathematics\par}}
\vspace{8pt}
{\color{rulegrey}\hrule height 0.6pt}
\vspace{7pt}
\thispagestyle{empty}

\textbf{Summary.} The root of the artifact is a file of five theorems, each proved in one line: the Riemann
hypothesis, in the formulation of the library and in the artifact's own; the non-vanishing of every Dirichlet
$L$-function on the half-plane \m{\re s>\tfrac12}; and the statement that the zeros of these functions in the
critical strip lie on the line \m{\re s=\tfrac12}, for all characters and for primitive ones. All five follow from
one theorem, the zero-free half-plane, by the functional equations. This note states the five theorems, shows how
each is obtained, and then says what one has to trust in order to believe them, and what that trust does not give.

\statement{leancol}{\fontsize{9.6}{13.2}\selectfont\textbf{What this note is.} The artifact \ArtifactName{} is a
body of Lean code: the compact edition of a formal proof, due to L.~Alpöge with the AI system Claude, whose final
theorems are stated below. That headline is a claim. The authors report that every file compiles with Lean v4.35.0
against a pinned version of the library Mathlib and that Lean's kernel accepted the proofs; an independent build is
recorded in the handout's record of checks when it reports. No mathematician has read the argument through, and
nothing here says that the Riemann hypothesis is proved. This note treats one unit of the artifact, its root: all
\UnitInNote{} of its \UnitMembers{} declarations. Their proofs are given in full; they are deductions of one line
from three results of another unit, which are stated as Inputs~1 to~3 and not proved here. Every statement was
compared with its Lean text. The Lean text read is \PinText; Appendix~A is the concordance of \ConcDecls{}
declarations.}

\sect{The five theorems}{s:intro}
Throughout, $\zeta$ is the Riemann zeta function and $L(s,\chi)$ the $L$-function of a Dirichlet
character~$\chi$, as the mathematical library of Lean defines them (§\ref{s:setting}).

\statement{arith}{\phantomsection\label{t:1}\textbf{Theorem 1} (the Riemann hypothesis, in the library's
formulation).\enspace \emph{For every complex number~$s$ with \m{\zeta(s)=0}, such that there is no natural
number~$n$ with \m{s=-2(n+1)}, and with \m{s\ne1}: \m{\re s=\tfrac12}.}

\vspace{4pt}\phantomsection\label{t:2}\textbf{Theorem 2} (the Riemann hypothesis, in the artifact's
formulation).\enspace \emph{For every complex number~$s$: if \m{\zeta(s)=0}, if there is no natural number~$n$ with
\m{s=-2(n+1)}, and if \m{s\ne1}, then \m{\re s=\tfrac12}.}}
\leanline{\lean{thmRH}\sep \lean{thmZeta}}

\statement{arith}{\phantomsection\label{t:3}\textbf{Theorem 3} (the zero-free half-plane).\enspace
\emph{For every integer \m{M\ge1}, every Dirichlet character~$\chi$ modulo~$M$ and every complex number~$s$ with
\m{\re s>\tfrac12}: if $\chi$ is not the principal character modulo~$M$ or \m{s\ne1}, then \m{L(s,\chi)\ne0}.}

\vspace{4pt}\phantomsection\label{t:4}\textbf{Theorem 4} (zeros in the critical strip).\enspace
\emph{For every integer \m{M\ge1}, every Dirichlet character~$\chi$ modulo~$M$ and every complex number~$s$: if
\m{0<\re s}, \m{\re s<1} and \m{L(s,\chi)=0}, then \m{\re s=\tfrac12}.}

\vspace{4pt}\phantomsection\label{t:5}\textbf{Theorem 5} (zeros in the critical strip, primitive
characters).\enspace \emph{For every integer \m{M\ge1} and every Dirichlet character~$\chi$ modulo~$M$ that is
primitive, and for every complex number~$s$: if \m{0<\re s}, \m{\re s<1} and \m{L(s,\chi)=0}, then
\m{\re s=\tfrac12}.}}
\leanline{\lean{thmHalf}\sep \lean{thmStrip}\sep \lean{thmPrim}}

Theorems 1 and~2 say the same thing: the first is stated with the proposition that the library itself calls the
Riemann hypothesis, the second with the artifact's copy of it. Theorem~4 is the generalised Riemann hypothesis for
Dirichlet $L$-functions in its usual form. Theorem~3 is the form in which the artifact proves it; it carries the
whole argument, and the other four are deduced from it in §\ref{s:proofs}.

\begin{center}
\begin{tikzpicture}[x=1mm,y=1mm,>={Stealth[length=1.6mm]},
  bx/.style={draw=rulegrey,line width=0.5pt,rounded corners=1.2pt,fill=arith!3!white,inner xsep=5pt,inner ysep=4pt,
    align=center,font=\fontsize{8.3}{10.4}\selectfont},
  gx/.style={bx,fill=inksoft!5!white},
  lb/.style={font=\fontsize{7.2}{8.8}\selectfont\color{inksoft},align=center}]
\node[gx] (abc) at (0,0) {Theorems A, B, C\\of the argument};
\node[gx] (star) at (32,0) {the main\\inequality};
\node[gx] (pr) at (60,0) {the bound for\\the primes};
\node[bx] (t3) at (90,0) {Theorem 3};
\node[bx] (t4) at (78,-19) {Theorem 4};
\node[bx] (t5) at (48,-19) {Theorem 5};
\node[bx] (t1) at (106,-19) {Theorems 1, 2};
\draw[->,inksoft] (abc) -- (star);
\draw[->,inksoft] (star) -- (pr);
\draw[->,inksoft] (pr) -- node[lb,above=1pt] {Landau} (t3);
\draw[->,inksoft] (t3) -- node[lb,left=2pt] {functional\\equation} (t4);
\draw[->,inksoft] (t3) -- node[lb,right=2pt] {\m{M=1}, functional\\equation} (t1);
\draw[->,inksoft] (t4) -- node[lb,above=1pt] {special case} (t5);
\end{tikzpicture}
\end{center}
\vspace{-4pt}
{\fontsize{8.9}{12}\selectfont\color{inksoft}\noindent The five theorems and what lies behind them. The grey boxes
and the two arrows marked with functional equations are the subject of the note \emph{The main inequality, the
bound for the primes, and the zero-free half-plane}; Theorems A, B and~C are the three main theorems of the argument
(short forms, enough forms, integral forms are long) and are read in later chapters.\par}

\sect{The setting}{s:setting}
\begin{conventions}\phantomsection\label{c:conventions}
The natural numbers are \m{0,1,2,\dots}, and $\re s$ is the real part of a complex number~$s$. The declarations of
this unit have no implicit arguments, so there is no ambient object to fix. The functions $\zeta$ and $L(s,\chi)$,
Dirichlet characters, principal and primitive characters are those of Mathlib, the mathematical library of Lean,
at the commit that the artifact pins. The note \emph{The four statements: the Riemann hypothesis and its
generalisation to Dirichlet $L$-functions} says what these objects are and what determines them. In short:
$\zeta$ and \m{L(\,\cdot\,,\chi)} are the analytic continuations of the series \m{\sum_{n\ge1}n^{-s}} and
\m{\sum_{n\ge1}\chi(n)n^{-s}}, with a conventional value at the pole \m{s=1} of $\zeta$ and of the $L$-function of
a principal character; a Dirichlet character modulo~$M$ is a multiplicative map from \m{\ZZ/M} to~$\CC$ that
vanishes outside the units, the principal character is the one with value~$1$ at every unit, and a character is
primitive when no character of a smaller modulus induces it.
\end{conventions}

\begin{definition}[the four statements]\label{d:statements}
The artifact defines four propositions.
\begin{enumerate}[label=\textup{(\arabic*)},leftmargin=2.3em,itemsep=2pt]
\item \kw{The Riemann hypothesis for $\zeta$}: for every complex number~$s$, if \m{\zeta(s)=0}, if there is no
natural number~$n$ with \m{s=-2(n+1)}, and if \m{s\ne1}, then \m{\re s=\tfrac12}.
\item \kw{The half-plane statement}: for every integer \m{M\ge1}, every Dirichlet character~$\chi$ modulo~$M$ and
every complex number~$s$ with \m{\re s>\tfrac12}, if $\chi$ is not the principal character modulo~$M$ or \m{s\ne1},
then \m{L(s,\chi)\ne0}.
\item \kw{The strip statement}: for every integer \m{M\ge1}, every Dirichlet character~$\chi$ modulo~$M$ and every
complex number~$s$, if \m{0<\re s}, \m{\re s<1} and \m{L(s,\chi)=0}, then \m{\re s=\tfrac12}.
\item \kw{The strip statement for primitive characters}: for every integer \m{M\ge1} and every primitive Dirichlet
character~$\chi$ modulo~$M$, and for every complex number~$s$, if \m{0<\re s}, \m{\re s<1} and \m{L(s,\chi)=0},
then \m{\re s=\tfrac12}.
\end{enumerate}
\end{definition}
\leanline{\lean{rh}\sep \lean{halfPlane}\sep \lean{strip}\sep \lean{stripPrim}\sep these four definitions form
another unit, the subject of the note \emph{The four statements}}

Theorems 2, 3, 4 and~5 assert propositions (1), (2), (3) and~(4), in this order. Theorem~1 asserts a proposition
of the library, whose text agrees with~(1) word for word, up to the way in which the number \m{n+1} is read as a
complex number.

\statement{inksoft}{\textbf{Inputs from another unit.} The three results below are taken for granted in this note.
They belong to the unit that ends the argument; they are proved in the note \emph{The main inequality, the bound
for the primes, and the zero-free half-plane}, the first from the main theorems of the argument, the other two
completely.

\vspace{4pt}\phantomsection\label{i:1}\textbf{Input 1} (the zero-free half-plane).\enspace
\emph{The half-plane statement holds.}

\vspace{3pt}\phantomsection\label{i:2}\textbf{Input 2} (from the half-plane to the strip).\enspace
\emph{If the half-plane statement holds, then the strip statement holds.}

\vspace{3pt}\phantomsection\label{i:3}\textbf{Input 3} (from the half-plane to $\zeta$).\enspace
\emph{If the half-plane statement holds, then the Riemann hypothesis for~$\zeta$ holds.}}
\leanline{\lean{grh}\sep \lean{zerosOnLine}\sep \lean{zetaZerosOnLine}}

Input~1 is the main theorem of the artifact. Inputs 2 and~3 are classical: the functional equation of the
completed $L$-function reflects a zero~$s$ with \m{0<\re s<\tfrac12} to the zero \m{1-s} of the $L$-function of the
inverse character, and for \m{M=1} the $L$-function is~$\zeta$.

\sect{Proofs}{s:proofs}
\begin{proof}[Proof of Theorem 3]
It is Input~1.
\end{proof}

\begin{proof}[Proof of Theorem 4]
By Input~1 the half-plane statement holds, and by Input~2 the strip statement follows.
\end{proof}

\begin{proof}[Proof of Theorems 1 and 2]
By Input~1 the half-plane statement holds, and by Input~3 the Riemann hypothesis for~$\zeta$ follows. This is
Theorem~2, and it is Theorem~1 because the proposition of the library and proposition~(1) of
Definition~\ref{d:statements} are the same.
\end{proof}

\begin{proof}[Proof of Theorem 5]
Let \m{M\ge1} be an integer and $\chi$ a primitive Dirichlet character modulo~$M$. Theorem~4, applied to $M$
and~$\chi$, gives the assertion for every complex number~$s$. That $\chi$ is primitive is not used.
\end{proof}
\leanline{each of the five formal proofs is one term\sep the terms use \lean{grh}, \lean{zerosOnLine} and
\lean{zetaZerosOnLine} exactly as above}

\sect{What one must trust}{s:trust}
The five theorems are short to state, and their formal proof is very long. What does a reader have to accept in
order to believe them, short of reading the proof? Three things, of different kinds; a fourth point says what
this trust does not give.

\subhead{1. That the statements mean what they say.} The four propositions of Definition~\ref{d:statements} are
written in a file that uses the library and nothing else, and the root file states the five theorems with these
propositions and with the library's own. So what is claimed can be read off two short files, and it rests on the
library's definitions of $\zeta$, of Dirichlet characters and of their $L$-functions. No definition made by the
artifact enters a statement: all its other definitions occur only inside proofs. The library's functions are
defined at every complex number, also at a pole, and the statements exclude the poles explicitly. The note
\emph{The four statements} goes through this.

\subhead{2. The library.} One trusts that the library's $\zeta$ and $L(s,\chi)$ are the classical functions.
This is not a matter of faith in a construction: the library proves that they agree with the Dirichlet series for
\m{\re s>1} and are holomorphic away from the pole, and these two properties determine them. What one accepts is
that these theorems of the library are themselves correctly checked, which is the next point.

\subhead{3. The kernel and the axioms.} A proof in Lean is a term that a small program, the kernel, checks
against the statement. The authors report three things: every file of the artifact compiles with the pinned
versions of Lean and of the library, without errors and without the placeholder that Lean accepts for a missing
proof; the five theorems depend on no axiom other than the three standard ones of Lean with classical logic
(propositional extensionality, the soundness of quotients, and the axiom of choice); and an independent tool,
which compares a file holding only the five statements with the root file and checks the proofs again, accepted
them. These are the authors' reports, read in the description that accompanies the artifact. On our side, a
reading of the text of all files of the artifact, made when this workspace was set up, found no placeholder and
no declaration of an axiom, and an independent build with the same versions is recorded in the handout's record of
checks when it reports. Which axioms a proof rests on is for the kernel to say, not for a reading of the text.
To accept the result on this ground is to trust the kernel of Lean, its implementation, and the three axioms.

\subhead{4. Nothing about the argument.} A kernel check establishes that the stated theorems follow from the
axioms. It does not tell a mathematician why they hold, and it is not a reading: at the time of writing no
mathematician has read the argument through. The purpose of this write-up is to make that reading possible. Each of its notes says
what was compared with the Lean text by a program, what was read by its writer, and what is taken for granted from
other units; the map of the handout shows which units have been read.

\sect{Notes on the formal text}{s:notes}
\begin{enumerate}[label=\textup{\arabic*.},leftmargin=1.6em,itemsep=5pt]
\item \emph{The five proofs are terms.} Theorem~3 is proved by the name of Input~1 alone: it is the same theorem
under a second name. Theorem~4 is Input~2 applied to Input~1. Theorem~1 is Input~3 applied to Input~1. Theorem~2
is proved by the name of Theorem~1. Theorem~5 takes the modulus, the assumption that it is not zero, the character
and the assumption that it is primitive, and returns Theorem~4 for that modulus and character.
\item \emph{Two formulations of the Riemann hypothesis.} Input~3 concludes proposition~(1) of
Definition~\ref{d:statements}, and Theorem~1 is stated with the proposition of the library; conversely Theorem~2 is
stated with~(1) and proved by Theorem~1. Lean accepts both steps only if the two propositions agree after
unfolding their definitions. So the acceptance of the file is also a check that the artifact's copy is the
library's formulation.
\item \emph{The hypothesis of Theorem~5.} The assumption that the character is primitive is received and not used.
Theorem~5 is stated because the generalised Riemann hypothesis is often formulated for primitive characters only;
for the zeros with \m{0<\re s<1} the two formulations are equivalent, as the note \emph{The four statements}
explains.
\item \emph{The modulus.} The formal text writes the modulus as a natural number~$N$ with the assumption
\m{N\ne0}, passed as an instance argument. Here it is an integer \m{M\ge1}; the letter~$N$ is kept in this
write-up for the level of the automorphic forms of the proof.
\item \emph{Authorship.} The root file carries one comment: the proof is Levent Alpöge's, with Claude. The
description of the artifact adds that the Lean text and the accompanying documents were written by Claude, and
that this edition is a shorter rewriting of an earlier formal proof by Alpöge.
\end{enumerate}

\subhead{How this note was written.} From the dossier of the unit. It has no structured proof, because each of
the five formal proofs is a single term; the program check of the dossier covered \UnitSteps{} leaf steps of
\UnitProofs{} proofs, found the five statements at their lines and compared the quotations of the imported
declarations with the source. I read the whole root file and the whole file of the four statements, and the
statements of the three inputs in their file. For §\ref{s:trust} I read the description that accompanies the
artifact, the file of five statements without proofs that it provides for the comparison tool, the configuration
of that tool, and the files that pin the versions of Lean and of the library; I also read the library's
formulation of the Riemann hypothesis in the library's public repository at the pinned commit. I did not run a
build. The diagram and §\ref{s:trust} are mine.

\head{Appendix A. Concordance with the Lean artifact}
{\fontsize{8.6}{11.6}\selectfont\renewcommand{\leanfont}{\ttfamily\fontsize{7.7}{10}\selectfont\color{leancol}}
\setlength{\LTleft}{0pt}\setlength{\LTright}{\fill}\setlength{\tabcolsep}{0pt}
\begin{longtable}{@{}p{30mm}p{56mm}p{43mm}r@{}}
\textbf{In this note} & \textbf{Lean declaration} & \textbf{File} & \textbf{Lines}\\\midrule\endhead
\ConcordanceRows
\end{longtable}}

\subhead{The five theorems as they are written in the artifact.}
\quoted{thmRH}
\quoted{thmZeta}
\quoted{thmHalf}
\quoted{thmStrip}
\quoted{thmPrim}
The three inputs:
\quoted{grh}
\quoted{zerosOnLine}
\quoted{zetaZerosOnLine}
In these texts {\leanfont RiemannHypothesis} is the proposition of the library, in the namespace of the library;
the four other propositions are those of Definition~\ref{d:statements}, and the table above gives their names in
the order (1) to~(4). The library's proposition is not quoted here, because the build of
this note copies quotations from the artifact only; it is to be read in the library, in the file of the Riemann
zeta function.

\head{Notation}
\begingroup\fontsize{8.6}{10.9}\selectfont\raggedright
\setlength{\columnsep}{9mm}\setlength{\parskip}{0pt}
\begin{multicols}{2}
\nt{$s$, $\re s$}{a complex number; its real part}
\nt{$n$}{a natural number, \m{n\ge0}}
\nt{$\zeta$}{the Riemann zeta function, as the library defines it}
\nt{$M$}{the modulus of a Dirichlet character, an integer \m{\ge1}}
\nt{$\chi$}{a Dirichlet character modulo~$M$}
\nt{$L(s,\chi)$}{the $L$-function of the Dirichlet character~$\chi$, as the library defines it}
\nt{$N$}{the modulus in the formal text; in this note it is~$M$}
\end{multicols}
\endgroup

\end{document}
